\documentclass[10pt,a4paper]{amsart}
\usepackage[margin=19mm]{geometry}
\usepackage{amsmath,amssymb,mathtools}
\usepackage[scr=rsfs,cal=euler]{mathalfa}
\usepackage{xfrac}
\usepackage[unicode,hidelinks,bookmarks=false]{hyperref}
\newcommand{\ie}{i.e.\ }
\newcommand{\cf}{cf.\ }
\newcommand{\resp}{respectively\ }
\newcommand{\Subsection}{\subsection}
\providecommand{\nobreakdash}{\nobreak\hbox{-}}
\setlength{\emergencystretch}{2em}
\multlinegap0pt
\usepackage{bm}
\usepackage{bbm}
\usepackage{dsfont}
\usepackage{xspace}
\usepackage{cancel}
\usepackage{mathtools}
\usepackage{amsthm}
\makeatletter
\@ifundefined{theorem}{\newtheorem{theorem}{Theorem}}{}
\makeatother
\title[$s$-Shunt Intersection Graph of a Graph]{$s$-Shunt Intersection Graph of a Graph}
\author{Vinny Susan Prebhath, Sudev Naduvath}
\date{Source: arXiv:2507.16309v1 (CC BY 4.0)}
\begin{document}
\maketitle

\noindent\emph{Source and licence.} $s$-Shunt Intersection Graph of a Graph. Version arXiv:2507.16309v1 (CC BY 4.0), distributed under the Creative Commons Attribution 4.0 International licence, as recorded in the arXiv licence metadata for this exact version and shown on the abstract page https://arxiv.org/abs/2507.16309v1. Full rights notice: licenses/arxiv-2507.16309-CC-BY-4.0.txt.\par\smallskip

\noindent\textit{Editorial scope.} Counted unit: the Theorem whose source label is \texttt{Asconn1} in the article (source0.tex lines 425--427, source label \texttt{Asconn1}), delivered together with the article's own complete original proof of it (source0.tex lines 428--445, one \texttt{proof} environment of 18 lines). No mathematics is written, repaired or re-derived: statement and proof are the article's own text line for line. Required context retained uncounted: none. Every definition, display and label of those lines is retained unchanged. Omitted from this package and not redistributed: 1--424, 446--751. Nothing inside a retained excerpt is omitted or altered: every retained body is delivered exactly as the article's lines are written, so no omission receipt of any kind accompanies this package. Citations: the retained excerpts carry no citation command at all, so no bibliography entry of the article is transcribed and no reference is redistributed. Reference adaptation: none was needed and none was made; every reference inside the delivered unit resolves inside it, so no unresolved reference or citation is produced. Printed slips kept literal: no printed word, symbol, hypothesis or number of the counted statement or of its proof was altered. Layout and class: the article's own class is \texttt{article} with packages including graphicx, etoolbox, amssymb, amsmath, amsthm, amsfonts, latexsym, correctmathalign, euscript, caption, subcaption, hyperref, geometry, authblk; the wrapper is the lane's pinned \texttt{amsart} editorial wrapper at 10pt on A4, which restates the theorem-like environments the retained text needs and adds the article's own macro definitions that the retained text needs (none), with extra wrapper packages (bm, bbm, dsfont, xspace, cancel, mathtools). The delivered numbering is the unit's own. Assets: no retained span contains a figure environment, a graphics-inclusion command, a PDF-page inclusion, a table or a TikZ picture; no image file is copied into this package, the package declares no asset dependency and no image file is redistributed with it. Licence: version arXiv:2507.16309v1 of this article is distributed under the Creative Commons Attribution 4.0 International licence, as recorded in the arXiv licence metadata for this exact version and shown on the abstract page https://arxiv.org/abs/2507.16309v1. A full rights notice is delivered at licenses/arxiv-2507.16309-CC-BY-4.0.txt. Characters: every retained excerpt is pure ASCII, so the delivered engine, XeLaTeX with its default Latin Modern text font, reports no missing character.
\begin{theorem} \label{Asconn1}
	The ssi-graph $A_{s}(G)$ of a graph $G$, if exists, is connected if and only if $G$ contains exactly one component with detour diameter at least $s+1$.
\end{theorem}

\begin{proof}
	%$G$ is connected here.
	\iffalse Assume that $A_{s}(G)$ is a non-empty disconnected graph. Hence, there are at least two components in $A_{s}(G)$. For each $s$-arc $\alpha_{i}$ in $G$, there exists at least one other $s$-arc $\alpha_{j}$ in $G$ such that $\alpha_{i}$ can be shunted onto $\alpha_{j}$. The $s$-arc $\alpha_{j}$ cannot be shunted onto $\alpha_{i}$. However, the $s$-arc obtained by reversing the order of the vertices of $\alpha_{j}$ can be shunted onto $\alpha_{i}'$ implying that the minimum number of vertices in each component of $A_{s}(G)$ is two. Since $A_{s}(G)$ is disconnected, the vertices in $G$, corresponding to the vertices of $A_{s}(G)$ in the first component, is not adjacent to any vertices in $G$, corresponding to the vertices of $A_{s}(G)$ in the other components. There may exist a path in $G$ between the vertices in $G$, corresponding to the vertices of $A_{s}(G)$ in the first component, and the vertices in $G$, corresponding to the vertices of $A_{s}(G)$ in the other components. This implies that either there exists an $s$-arc, that is a vertex of $A_{s}(G)$, with vertices from an $s$-arc in the first component of $A_{s}(G)$ and vertices from an $s$-arc in any other component of $A_{s}(G)$, implying that the number of components of $A_{s}(G)$ is less, or $G$ contains at least two components with detour diameter at least $s+1$ resulting in a contradiction. Therefore, $A_{s}(G)$ is connected. \fi
	
	%If $A_{s}(G)$ is connected, then there exists a path in $G$ between the end vertices of any two $s$-arcs of $G$. If there exists no path between the end vertices of these two $s$-arcs in $G$, then the corresponding $s$-arcs are not vertices in $A_{s}(G)$ because this implies that these $s$-arcs cannot be shunted onto any other $s$-arc in $G$. Hence, $G$ is graph with exactly one component whose detour diameter is at least $s+1$.
	
	%Suppose $G$ is a disconnected graph with at least two components whose detour diameter is at least $s+1$. 
	
	%Required: If $G$ has more than one such components, then $A_{s}(G)$ is disconnected implying that if $A_{s}(G)$ is connected, then it is necessary that $G$ must have exactly one such component. Conversely, if $G$ has exactly one component with the condition, then $A_{s}(G)$ is connected.
	
	%We first show that the ssi-graph of a connected graph $G$ with detour diameter at least $s+1$ is always connected. 
	Let $G$ be a connected graph with detour diameter at least $s+1$. Assume that the ssi-graph $A_{s}(G)$ is disconnected. For each $s$-arc $\alpha_{i}$ in $G$ that is a vertex of $A_{s}(G)$, there exists at least one other $s$-arc $\alpha_{j}$ in $G$ such that $\alpha_{i}$ can be shunted onto $\alpha_{j}$. Also, the $s$-arc $\alpha_{j}'$ can be shunted onto $\alpha_{i}'$ implying that each component of $A_{s}(G)$ consists of at least two vertices. Since $A_{s}(G)$ is disconnected, the vertices of the $s$-arc, say $\alpha$ of $G$, corresponding to the vertex $\alpha$ in a component $Q_{i}$ of $A_{s}(G)$, are not adjacent to any vertices of the $s$-arc, say $\beta$ of $G$, corresponding to the vertex $\beta$ in any other component $Q_{j};i\ne j$ of $A_{s}(G)$. However, there exists at least one path $P$ in $G$ with one end vertex in $\alpha$ and the other in $\beta$. Then, there exists a shuntable $s$-arc $\gamma$ in $G$ consisting of vertices of $\alpha,\beta$ and $P$. Hence, $\gamma$ is a vertex of $A_{s}(G)$ which is adjacent to both $\alpha$ and $\beta$, contradicting the fact that $\alpha$ and $\beta$ are in different components of $A_{s}(G)$. Therefore, the ssi-graph of a connected graph is always connected. Also, note that any component of $G$ having detour diameter less than $s+1$, if exists, does not contribute any vertex to $A_{s}(G)$. %the vertex set of the ssi-graph of graphs with detour  diameter less than $s+1$ is always empty. Therefore, if $G$ has more than one such components, each component will yield a connected component in $A_{s}(G)$. But since the vertex set of $s$-arcs from the disjoint components are distinct, the resulting ssi-graph is disconnected. Hence, the number of components of $A_{s}(G)$ is same as the number of components of $G$ with detour diameter at least $s+1$. Hence, if $G$ is a graph with exactly one component with detour diameter at least $s+1$, then $A_{s}(G)$ is connected.
	
	%Since $A_{s}(G)$ is disconnected, the vertices of $s$-arcs of $G$, corresponding to the vertices of a component of $A_{s}(G)$, are not adjacent to any vertices of the $s$-arcs of $G$, corresponding to the vertices of any other components of $A_{s}(G)$. However, there exists a path in $G$ between the vertices in $G$, corresponding to the vertices of $A_{s}(G)$ in the first component, and the vertices in $G$, corresponding to the vertices of $A_{s}(G)$ in the other components. This implies that there exists an $s$-arc, that is a vertex of $A_{s}(G)$, with vertices in $G$ from an $s$-arc in the first component of $A_{s}(G)$ and vertices in $G$ from an $s$-arc in some other component of $A_{s}(G)$, implying that the number of components of $A_{s}(G)$ is less than $k$, a contradiction. Therefore, the ssi-graph of a connected graph is always connected. Also, the vertex set of the ssi-graph of graphs with detour  diameter less than $s+1$ is always empty. Therefore, if $G$ has more than one such components, each component will yield a connected ssi-graph. But since the vertex set of $s$-arcs from the disjoint components are distinct, the resulting ssi-graph is disconnected. Hence, the number of components of $A_{s}(G)$ is same as the number of components of $G$ with detour diameter at least $s+1$. Hence, if $G$ is a graph with exactly one component with detour diameter at least $s+1$, then $A_{s}(G)$ is connected. %then it is necessary that $G$ must contain at most one component with detour diameter at least $s+1$.
	
	If $A_{s}(G)$ is a non-empty connected graph, then there exists a path in $G$ between the end vertices of any two $s$-arcs of $G$, that can be shunted onto some other $s$-arc of $G$. For a graph $G$ with detour diameter at least $s+1$, every vertex of $G$ need not be a vertex of any $s$-arc of $G$, that can be shunted onto some other $s$-arc of $G$. If there exists a vertex $v$ of $G$ that is not a vertex of any of the $s$-arcs of $G$, then the detour eccentricity of $v$ is less than $s+1$. As a consequence, there may exist a component in $G$ with detour diameter less than $s+1$. %Hence, these vertices with detour eccentricity less than $s+1$ are vertices of a component of $G$ with either detour diameter at least $s+1$ or detour diameter less than $s+1$. 
	Therefore, all vertices of $A_{s}(G)$ must be $s$-arcs from the same component of $G$ with detour diameter at least $s+1$.
\end{proof}

\end{document}
